\chapter{حجية خبر الآحاد وشروط الحديث الصحيح}

\section{مقدمة في خبر الخاصة (خبر الآحاد)}
الحمد لله والصلاة والسلام على رسول الله، أما بعد؛ فقد شرع الإمام الشافعي رحمه الله تعالى في بيان باب عظيم من أبواب أصول الفقه والحديث، وهو "باب خبر الواحد".
وقد سأله سائل: حدد لي أقل ما تقوم به الحجة على أهل العلم حتى يثبت عليهم خبر الخاصة؟ 
و(خبر الخاصة) عند الشافعي يقابله (خبر العامة)؛ فخبر العامة هو المتواتر الذي يعلمه عامة المسلمين (كفرضية الصلوات الخمس)، أما خبر الخاصة فهو ما ينقله ويعلمه أهل العلم وطلابه، وهو ما نطلق عليه اصطلاحاً (خبر الآحاد).
فأجاب الشافعي: أقل ما تقوم به الحجة هو خبر الواحد عن الواحد حتى يُنتهى به إلى النبي ﷺ، أو من يُنتهى به إليه دونه من الصحابة والتابعين.

\section{شروط الحديث الصحيح عند الإمام الشافعي}
قرر الشافعي رحمه الله أنه لا تقوم الحجة بخبر الخاصة (الآحاد) إلا باجتماع أمور وشروط في رواته، وهي التي استقر عليها علماء مصطلح الحديث في تعريف "الحديث الصحيح"، وهي:
\begin{enumerate}
    \item \textbf{العدالة:} أن يكون الراوي ثقة في دينه، معروفاً بالصدق في حديثه (العدل).
    \item \textbf{الضبط والفهم:} أن يكون عاقلاً فاهماً لما يحدّث به، عالماً بما يحيل معاني الحديث من الألفاظ (إذا كان يروي بالمعنى).
    \item \textbf{الأداء كما سمع:} أن يكون ممن يؤدي الحديث بحروفه كما سمع إذا لم يكن عالماً بلغة العرب وما يحيل المعاني؛ لئلا يحوّل الحلال إلى حرام بخطأ في اللفظ.
    \item \textbf{ضبط الصدر أو الكتاب:} أن يكون حافظاً إن حدّث من حفظه، وضابطاً لكتابه حافظاً له إن حدّث من كتابه.
    \item \textbf{عدم الشذوذ (موافقة الثقات):} إذا شَرِك أهل الحفظ في الحديث وافق حديثهم، ولا يخالف الثقات.
    \item \textbf{السلامة من التدليس:} أن يكون بريئاً من أن يكون مدلّساً؛ يحدّث عمن لقي ما لم يسمع منه.
    \item \textbf{اتصال السند:} أن يكون السند موصولاً إلى النبي ﷺ أو إلى من دونه، فكل راوٍ في السند يُشترط فيه هذه الشروط ليكون مثبتاً لمن حدّثه وعمن حدّث عنه.
\end{enumerate}

\section{الفرق بين الشهادة والرواية}
طلب السائل من الشافعي أن يضرب له مثلاً ويقيس خبر الآحاد على شيء من الشهادات التي يعرفها الناس. فبين الشافعي أن الرواية والشهادة يجتمعان في أشياء (كاشتراط العدالة والصدق)، لكنهما يفترقان في أمور جوهرية:
\begin{itemize}
    \item \textbf{في العدد والذكورة:} تُقبل في الرواية رواية الواحد ورواية المرأة الواحدة، بينما لا تُقبل في الشهادة شهادة الواحد وحده ولا المرأة وحدها.
    \item \textbf{في صيغة الأداء:} يُقبل في الرواية صيغة (عنعنة) الثقة غير المدلس كقوله: "حدثني فلان عن فلان"، بينما لا تُقبل في الشهادة إلا بالتصريح القطعي: "سمعتُ، أو رأيتُ، أو أشهدني".
    \item \textbf{في الترجيح:} تختلف الأحاديث فيُرجَّح ببعضها استدلالاً بكتاب أو سنة أو إجماع أو قياس، وهذا لا يوجد في باب الشهادات.
\end{itemize}

\subsection{لماذا يُحتاط للرواية أكثر من الشهادة؟}
بيّن الشافعي أن إحالة معنى الحديث أخفى من إحالة معنى الشهادة. فالشاهد يُنشئ كلاماً يعبر به عما رأى بلفظه، والخطأ في الشهادة يضر فرداً بعينه (في مال أو عقوبة). أما المحدّث فهو ناقل لحكم شرعي يحلل ويحرم للأمة جمعاء، والخطأ في لفظة قد يحيل الحلال حراماً؛ لذا كان الاحتياط في الرواية والتشديد في شروط رواتها (كالضبط والفهم والسلامة من التدليس) أشد من الاحتياط في الشهود.

\section{حكم رواية المدلّس وتعديل الرواة}
قرر الشافعي مذهبه الصارم في التدليس قائلاً: «ومن عَرَفناه دلَّس مَرَّة، فقد أبان لنا عورته في روايته». فمن دلّس (بأن روى عمن لقيه ما لم يسمع منه بصيغة موهمة كـ "عن")، لا يُقبل حديثه حتى يصرّح بالتحديث (سمعت أو حدثني). والتدليس ليس كذباً محضاً، إذ لو كان كذباً لردت روايته بالكلية، لكنه أظهر عورة في دقة نقله، فوجب التوقف في عنعنته.
كما بيّن أن رواية الثقة عن شيخ لا نعرفه لا تُعد توثيقاً له، بل لا بد من تعديل صريح للراوي.

\section{التحذير من الكذب على رسول الله ﷺ}
ختم الشافعي هذا المبحث ببيان ورع المحدثين وحرصهم على ألا يكذبوا على رسول الله ﷺ؛ مستدلاً بالأحاديث المتواترة في ذلك كقوله ﷺ: «من يقل عليّ ما لم أقل فليتبوأ مقعده من النار»، وقوله: «حدثوا عن بني إسرائيل ولا حرج، وحدثوا عني ولا تكذبوا عليّ».

فبيّن أن النبي ﷺ لم يُبح الكذب في أخبار بني إسرائيل، بل أباح التحديث عنهم على جهة البلاغ لتعذر الأسانيد إليهم ولطول المدة، ولكنه شدّد في النقل عنه ﷺ بألا يكون إلا بنقل الإسناد والتثبت الدقيق، والاحتراز التام من الكذب عليه، فكلامه دين وشرع إلى يوم القيامة، ولا يستدل على الصدق والكذب إلا بمعرفة الرواة ونقل الإسناد.